\par\Needspace{8\baselineskip}\bigskip
\SourceLeaf{457}
\pdfbookmark[3]{Exposition, 9}{pallas-II-article9}
\begin{center}9.\end{center}
\noindent\hspace*{\parindent}Let us suppose that the perturbations add the quantities \(\xi,\eta,\zeta\) to the fundamental equations, so that we have
\[
\begin{aligned}
0&=\dfrac{\mathrm{dd}x}{\mathrm{d}t^2}+\dfrac{mx}{r^3}+\xi\\[5pt]
0&=\dfrac{\mathrm{dd}y}{\mathrm{d}t^2}+\dfrac{my}{r^3}+\eta\\[5pt]
0&=\dfrac{\mathrm{dd}z}{\mathrm{d}t^2}+\dfrac{mz}{r^3}+\zeta.
\end{aligned}
\]
\par We see that \(\xi,\eta,\zeta\) are the forces obtained by resolving, along directions opposite to those of the three axes, all the forces that perturb the planet's motion relative to the Sun. We can also imagine these forces combined into a single force whose magnitude is \(=\chi\) and whose direction is parallel to a radius of our sphere drawn from the center toward the point \(Q\) on the surface. We then have
\[
\begin{aligned}
\xi&=-\chi\cos XQ\\
\eta&=-\chi\cos YQ\\
\zeta&=-\chi\cos ZQ.
\end{aligned}
\]
\par In place of the forces \(\xi,\eta,\zeta\), it may be useful to introduce those obtained by resolving \(\chi\) along three other mutually perpendicular directions. Let us first suppose that the first acts perpendicular to the plane of the osculating orbit, the second in the direction of the radius vector, and consequently the third perpendicular to the radius vector in the plane of the orbit, or, more precisely, that these three components act parallel to the three radii of the sphere drawn toward the points \(M,L,\Lambda\), so that they become
\[
\begin{aligned}
\chi\cos MQ&=-\xi\cos XM-\eta\cos YM-\zeta\cos ZM\\
\chi\cos LQ&=-\xi\cos XL-\eta\cos YL-\zeta\cos ZL\\
\chi\cos\Lambda Q&=-\xi\cos X\Lambda-\eta\cos Y\Lambda-\zeta\cos Z\Lambda.
\end{aligned}
\]
\par Thus the first will be
\[
=-\dfrac{A\xi+B\eta+C\zeta}{\BareRoot{mp}}\mathpunct{\raisebox{.45ex}{,}}
\]
\SourceLeaf{458}the second
\[
=-\dfrac{x\xi+y\eta+z\zeta}{r}\mathpunct{\raisebox{.45ex}{,}}
\]
the third
\[
=-\dfrac{(Bz-Cy)\xi+(Cx-Az)\eta+(Ay-Bx)\zeta}{r\BareRoot{mp}}\mathpunct{\cdot}
\]
We shall denote these three forces by \(mW,mT,mV\).
